\section{Conclusion}\label{sec:conclusion}

For unequal couplings the partition ratio itself can be bounded:
Theorem~\ref{cl:thm:unequal-gap} gives a physical lattice gap uniform
in spatial volume. The continuous-time $\SU(2)$ result also retains
spectral weight at a finite energy in a gauge-invariant observable.
Extending that proof to the isotropic tuning would require controlling
the spatial plaquette action when its coefficient grows along with the
temporal one; the absolute expansion in Section~\ref{cl:sec:unequal}
loses smallness in precisely that regime.

The strengthened criterion of Theorem~\ref{sc:thm:continuum} replaces the
window of tori by one cube: the inequality
$\psi_{(2T_0)^3}(T_0)\le1/4097$ at $T_0\le Nt_0$ implies
$\Delta=t_0^{-1}\log(63/32)$. The proof derives vacuum dominance and the
lattice-to-physical spectral conversion directly from that datum. The
zero--one law of Theorem~\ref{sc:thm:zero-one} further shows that the
normalised physical cube pressure either stays above a fixed threshold
at every scale or enters the exponentially decaying branch. It does not
select the branch of the tuned Wilson system.

Theorem~\ref{cl:thm:rough-cells} controls joint rough-cell events at weak
coupling. On even cubes the removal of all rough clusters of size at
least half the time period changes the seam probability by the explicit
error in Corollary~\ref{cl:cor:seam-clusters}. It remains to compare the
weights of the two matchings on the retained configurations. On the cube
of side $24T$, Theorem~\ref{cl:thm:seam-aspect} reduces the sufficient
unconditional seam probability to $97/266$. Proposition~\ref{cl:prop:seam-winding}
shows why the comparison has to average over the spatial fields.
For $\SU(2)$, Proposition~\ref{cl:prop:row-pins} permits all straight
windings in one spatial direction to be fixed to the identity with a
vanishing error in that comparison when $n/(1+4b)-3\log n\to+\infty$.
The remaining integrals couple the transverse layers at the original
Wilson coupling. Theorem~\ref{cl:thm:charged-gap} bounds the charged
spectrum uniformly in volume for every compact connected matrix group;
the gauge-invariant excited spectrum remains the physical target.
At $T\asymp N$ and $b=O(\log N)$ the same seam odds can be computed
before gauge projection with error smaller than every inverse power of
$N$ (Corollary~\ref{cl:cor:temporal-physical}).

The replica construction gives the exact single-cube target
$I_{b,(2T)^3,T}\le1/(16388T^3)$. Its conditional trace counterpart
requires bounds in squared-likelihood tilted histories, rather than the
stationary average alone. The large-field witness rules out a uniform
global diameter certificate, and the Wilson Dirichlet identity isolates
the upper energy and lower variance estimates needed to retain a
continuum excitation. None of these identities supplies the corresponding
uniform estimate for the tuned Wilson measure.


The conditional Wilson comparison of Theorem~\ref{cl:thm:link-poincare}
is independent of volume and includes every prefix history. Its constant
grows as $e^{12b}$, so it does not supply a coupling-uniform gap.
Corollary~\ref{cl:cor:wilson-score} identifies the remaining score budget
that would suffice for the single-cube datum. The bounded-energy
criterion of Proposition~\ref{cl:prop:spectral-mass} permits retained
observables with a divergent high-energy component, and
Appendix~\ref{cl:app:flow} gives a regularisation that preserves
positive-time support. Neither construction proves the required
Wilson spectral-mass lower bound or the continuum convergence estimates.

The full four-dimensional weak-coupling problem remains unresolved:
the datum must still be proved along the continuum tuning, and a
nontrivial continuum theory satisfying the required axioms must be
constructed for every compact simple group. The single-cube criterion
complements the following strong-coupling and slab results.



Lattice Yang--Mills theory in four dimensions has a mass gap with the explicit
lower bound $m_G(\beta)$ for every compact simple gauge group in the range
$N\beta<\frac2{29}\kappa_G$, and for $\SU(2)$ up to $\beta_W=\frac13$ with
the bound $-\log(3\beta_W)$ and up to $\beta_W=\frac9{25}$ by resampling stars. A
single
spectral gap $\varpi$ of the dynamics that resamples cubes of a fixed physical
size $\ell$, uniform along the continuum tuning, gives every nontrivial
continuum limit the mass gap $16\varpi/(625e^2\ell)$. That spectral gap holds
at strong coupling, for $\SU(2)$ up to $\beta_W=\frac7{20}$ uniformly in the
cube size and for $\SU(3)$ up to $\beta_W=0.45$, and for massive free fields;
it fails for massless fields, at fixed cube size at weak coupling, and in the
abelian Coulomb phase.

For a continuum theory built in the other order, by scaling limits on tori
of fixed physical side, the same cube condition on cylinders, together with
the independence of the volume limit from the order of its two limits, gives
the gap $16\varpi/(625e^2\ell)$, a unique vacuum, and a single phase
(Theorem~\ref{thm:E}); the proof passes through the transfer matrix of the
cylinder and the Euclidean vectors of the Wightman theory, and uses nothing
beyond the three hypotheses. The independence of the order of the limits
follows from a single renormalisation-free statement, that the partition
function of the symmetric torus relative to the vacuum of the cylinder tends
to one (Corollary~\ref{cor:Z}), and the cube condition on cylinders follows
from the slab decorrelation of Theorem~\ref{thm:D}, which holds on cylinders
(Proposition~\ref{prop:bridgeM}). With a construction of the torus theories
taken as the existence input, the statement of the Clay problem for those
theories reduces to these two conditions; neither is proved here.

All of these conditions follow from one. If, at a single time extent of
fixed physical size, the thermal part of the free energy of the tori whose
sides are between two and forty-eight times that extent is at most a small
explicit multiple of their volume, then the ascent of
Theorem~\ref{thm:ascentfree} gives the gap on every larger torus and the
vacuum dominance~(Z), and Theorem~F gives the Wightman theory the mass gap
$0.172/t_0$ and a unique vacuum. The statement of the Clay problem for the
theories of~(W) is thereby reduced to (W) itself and to this one bound,
which is an inequality between two partition functions. The bound contains
the lattice gap at weak coupling and is open; and the constructions that
would avoid it, by taking the transporters without renormalisation, by
lifting the two-dimensional theory, or by keeping the lattice, give trivial
theories or violate the axioms (Appendix~\ref{app:degenerate}).

The cube gap follows from a decay of correlations, below the sharp threshold
$\frac12$, between the gauge-invariant observables on the two sides of a slab
of physical width $\ell/2$, uniformly in the boundary condition. This
hypothesis is a form of mass generation at one scale: it is a bound on the
maximal correlation of two layers of a box, a quantity which in the vacuum
of a cylinder equals exactly $e^{-m_{\mathrm{lat}}s}$ at distance $s$
(Proposition~\ref{prop:vacuumlayers}), so that its whole content is the
uniformity in the boundary condition. It holds in two dimensions along the
weak-coupling tuning and fails there for strips one link wide; in four
dimensions it is not proved. Curvature arguments, with or
without gauge fixing, stop at explicit couplings of order one; spectral gaps
for blocks of fixed lattice size force a trivial continuum limit; and a
lattice gap tending to zero at the right rate produces a gap in a continuum
limit that is not identified with Yang--Mills theory. Statements (D) and (E)
remain open, as do (U), (Z) and (D$_\ell$) for the theories built on tori,
and with them the Clay problem; (Z) holds at strong coupling and in two
dimensions and follows from an entropy bound on the cylinder, and no
correlation, gap, slab or cube hypothesis reaches it
(Section~\ref{ssec:thermalvacuum}, Proposition~\ref{prop:Ztwod}), and it
reduces exactly to an entropy bound at one physical scale together with a
Casimir-energy condition on the halving of the spatial torus
(Proposition~\ref{prop:Zonescale}), both of which are doubling defects of
the free energies of finite tori, out of reach of the gap and of the free
energy per site, and the second of which follows from a finite-size bound
of massive type (Propositions~\ref{prop:Zdoubling}, \ref{prop:ringsfour}
and~\ref{prop:FSC}); given the gap on the tori of the halving chain, the
Casimir condition reduces to the physical Casimir energy lying below the
gap, and (C) to the exclusion of a linearly growing branch of thermal
corrections (Theorem~\ref{thm:gapcasimir}, Proposition~\ref{prop:dichotomy}),
a branch that positivity and the gap exclude once the one-scale bound is
taken on elongated boxes, so that (Z) then follows with no hypothesis at a
second scale (Theorem~\ref{thm:ascentZ}, Corollary~\ref{cor:branch}),
and the free field shows that (C) is a statement about tori of about eighteen
correlation lengths, violated by the free gluons at every scale
(Corollary~\ref{cor:freeC}, Proposition~\ref{prop:freegluons}); while the
slab hypothesis is a block Poincar\'e inequality uniform in the boundary
condition and the norm of a conditional expectation between the two sides
(Section~\ref{ssec:blockgap}, Proposition~\ref{prop:rhonorm}), and it is
not a consequence of a gap of the transfer matrix
(Propositions~\ref{prop:fieldrings} and~\ref{prop:fieldringsall}), while
the patch inequality that implies it is exactly the local Knabe inequality
on the torus (Proposition~\ref{prop:fibre}); what this paper contributes is an
explicit strong-coupling theory for every gauge group, a reduction of the
continuum gap to one estimate on one scale and of that estimate to a decay of
correlations across one slab, and a map of what the present methods can reach.
